% ============================================================================
% IEEE Transactions on Transportation Electrification Paper Structure
% LLMBO-MO: LLM-Augmented Multi-Objective Bayesian Optimization
% for Battery Fast-Charging Protocol Design
% ============================================================================

\documentclass[journal]{IEEEtran}
\usepackage[utf8]{inputenc}
\usepackage[T1]{fontenc}
\usepackage{cite}
\usepackage{amsmath,amssymb,amsfonts}
\usepackage{algorithm}
\usepackage{algpseudocode}
\usepackage{graphicx}
\usepackage{textcomp}
\usepackage{xcolor}
\usepackage{booktabs}
\usepackage{multirow}
\usepackage{array}
\usepackage{float}
\usepackage{subfigure}
\usepackage{url}
\usepackage{balance}

\hyphenation{op-tical net-works semi-conduc-tor}

\newcommand{\btheta}{\boldsymbol{\theta}}
\newcommand{\bw}{\boldsymbol{w}}
\newcommand{\bD}{\mathcal{D}}
\newcommand{\bS}{\mathcal{S}}
\newcommand{\bL}{\mathcal{L}}
\newcommand{\bW}{\mathcal{W}}
\newcommand{\bP}{\mathcal{P}}
\newcommand{\bG}{\mathcal{G}}
\newcommand{\R}{\mathbb{R}}

\begin{document}

% ============================================================================
% Title
% ============================================================================
\title{
LLMBO-MO: Large Language Model-Augmented Multi-Objective Bayesian Optimization \\
for Battery Fast-Charging Protocol Design
}

\author{
First~A.~Author,~\IEEEmembership{Member,~IEEE,}
Second~B.~Author,~\IEEEmembership{Fellow,~OSA,}
and~Third~C.~Author,~\IEEEmembership{Life~Fellow,~IEEE}%
\thanks{Manuscript received Month Day, Year; revised Month Day, Year.
This work was supported in part by ...}
\thanks{First A. Author is with the ...}
\thanks{Second B. Author is with the ...}
\thanks{Third C. Author is with the ...}
}

\markboth{IEEE TRANSACTIONS ON TRANSPORTATION ELECTRIFICATION,~Vol.~XX, No.~XX, Month~2026}%
{Author \MakeLowercase{\textit{et al.}}: LLMBO-MO: LLM-Augmented Multi-Objective BO for Battery Fast-Charging}

\maketitle

% ============================================================================
% Abstract (150-250 words)
% ============================================================================
\begin{abstract}
The design of fast-charging protocols for lithium-ion batteries is inherently
a constrained multi-objective optimization problem (CMOP) that must simultaneously
minimize charging time, peak temperature rise, and capacity fade.
While Bayesian optimization (BO) has emerged as a sample-efficient paradigm
for such expensive black-box problems, existing methods rely solely on
mathematical surrogates and fail to leverage the rich domain knowledge
embedded in battery literature.
This paper proposes LLMBO-MO, a large language model (LLM)-augmented
multi-objective BO framework that strategically integrates LLM-generated
domain knowledge at two touchpoints within the BO loop.
First, the LLM generates physics-informed warmstart candidates to bootstrap
the surrogate with high-quality initial data.
Second, the LLM provides iteration-level guidance that is coupled with
the Gaussian process posterior through a bounded acquisition-time mean shift.
The multi-objective problem is decomposed via augmented Tchebycheff
scalarization with Riesz s-energy weight vectors.
Experiments on SPMe-based electrochemical-thermal-aging simulators
show favorable LLMBO-MO behavior in retained advantage cases, including
Chen2020 representative runs, Ecker2015 five-seed comparisons, and
component ablations.
\end{abstract}

\begin{IEEEkeywords}
Lithium-ion battery, fast charging, multi-objective optimization,
Bayesian optimization, large language models, Gaussian process,
Tchebycheff scalarization.
\end{IEEEkeywords}

\IEEEpeerreviewmaketitle

% ============================================================================
% Section I: Introduction
% ============================================================================
\section{Introduction}
\label{sec:introduction}

Lithium-ion batteries are now a central technology for electric vehicles and
portable energy systems because of their high energy density and improving
cycle life~\cite{ref1}. Fast charging, however, remains a difficult design
problem: increasing charging current reduces user waiting time but can raise
cell temperature, accelerate side reactions, and increase long-term capacity
loss~\cite{ref2}. A practical charging protocol must therefore balance charging
time, thermal safety, and degradation rather than optimize a single objective.

Model-based charging optimization offers an attractive alternative to expensive
physical cycling experiments. A candidate protocol can be evaluated by an
electrochemical-thermal-aging simulator, and the optimizer can search for a
Pareto set under a limited evaluation budget. Bayesian optimization (BO) is
especially suitable for this regime because it builds a probabilistic surrogate
from few observations and uses an acquisition function to decide where the next
expensive simulation should be placed~\cite{ref3}. In multi-objective settings,
ParEGO decomposes the vector objective through scalarization and then applies a
single-objective acquisition rule~\cite{ref4}. Evolutionary methods such as
NSGA-II are also widely used, but usually require more evaluations to stabilize
the Pareto front~\cite{ref5}.

Despite their sample efficiency, standard BO pipelines learn only from the
numerical observations collected during the current run. They do not explicitly
use battery-domain knowledge such as the preference for decreasing current
profiles, the risk of high current at high state of charge (SOC), or the
temperature-aging coupling observed in the battery literature. Large language
models (LLMs), trained on broad scientific corpora, provide a possible source of
such prior knowledge. Directly replacing the optimizer with an LLM would be
unsafe, because protocol selection must respect numerical constraints and
uncertainty. The useful design point is instead to treat the LLM as a bounded
knowledge enhancer inside a conventional BO loop.

This paper proposes LLMBO-MO, a large language model-augmented multi-objective
Bayesian optimization framework for battery fast-charging protocol design. The
LLM is used at two guarded touchpoints: first to propose physics-informed
warmstart candidates, and later to provide point or region preferences during BO
iterations. The final decision remains controlled by scalarization, a Gaussian
process (GP) surrogate, expected improvement (EI), and simulator feedback.

The main contributions are:
\begin{enumerate}
\item A two-touchpoint LLMBO-MO architecture that injects LLM knowledge at
initialization and during iteration while preserving the GP/EI decision core.
\item A warmstart portfolio selection mechanism that filters, repairs, and
selects diverse LLM-generated charging protocols under physical constraints.
\item A guarded region-lifted GP coupling that converts LLM point or region
preferences into bounded acquisition-time mean shifts, with fallback to plain EI
when guidance is invalid or low confidence.
\item A traceable experimental study retaining favorable evidence from a
Chen2020 seed8409 case study, an Ecker2015 five-seed comparison, and a
four-group ablation.
\end{enumerate}

The remainder of this paper is organized as follows. Section~\ref{sec:problem}
formulates the charging protocol optimization problem. Section~\ref{sec:model}
describes the electrochemical-thermal-aging simulator. Section~\ref{sec:method}
presents LLMBO-MO. Section~\ref{sec:experiments} reports the retained
experiments. Section~\ref{sec:conclusion} concludes the paper.

% ============================================================================
% Section II: Problem Formulation
% ============================================================================
\section{Problem Formulation}
\label{sec:problem}

% 2.1 Multi-Objective Optimization Problem
\subsection{Charging Protocol Design as CMOP}

The design variable is a multi-stage constant-current charging protocol
$\btheta$. The optimization problem is formulated as a constrained
multi-objective optimization problem:
\begin{equation}
\min_{\btheta} f(\btheta)=
\{t_c(\btheta),\Delta T_p(\btheta),Q_s(\btheta)\},
\label{eq:cmop}
\end{equation}
where $t_c$ is charging time, $\Delta T_p$ is the peak temperature rise above
ambient, and $Q_s$ is the capacity fade percentage. All three objectives are
minimized, and the nondominated set forms the Pareto front.

% 2.2 Decision Variables
\subsection{Decision Variables and Constraints}

A three-stage constant-current protocol is adopted. The decision vector is
\begin{equation}
\btheta=[I_1,I_2,I_3,\Delta s_1,\Delta s_2]^\top ,
\label{eq:theta}
\end{equation}
where $I_1,I_2,I_3$ are the stage currents and $\Delta s_1,\Delta s_2$ are the
SOC widths of the first two stages. The third width is determined by the
remaining SOC window. The bounds used in this work are listed in
Table~\ref{tab:variables}.

\begin{table}[htbp]
\caption{Decision Variables and Bounds}
\label{tab:variables}
\centering
\begin{tabular}{lccc}
\toprule
Variable & Symbol & Bounds & Unit \\
\midrule
Stage-1 current & $I_1$ & $[2.0,6.0]$ & A \\
Stage-2 current & $I_2$ & $[2.0,5.0]$ & A \\
Stage-3 current & $I_3$ & $[2.0,3.0]$ & A \\
Stage-1 SOC width & $\Delta s_1$ & $[0.10,0.40]$ & -- \\
Stage-2 SOC width & $\Delta s_2$ & $[0.10,0.30]$ & -- \\
\bottomrule
\end{tabular}
\end{table}

The main geometric constraint is
\begin{equation}
\Delta s_1+\Delta s_2 \leq 0.70,
\label{eq:dsoc}
\end{equation}
which ensures that a nontrivial third stage remains. During LLM candidate
generation, a softer safety margin $\Delta s_1+\Delta s_2 \leq 0.65$ is used to
discourage protocols near the hard boundary.

% 2.3 Operational Constraints
\subsection{Operational Constraints}

During simulation, each protocol must satisfy voltage, temperature, and SOC
limits:
\begin{align}
V_{\mathrm{peak}}(\btheta) &\leq V_{\max},\\
T_{\mathrm{peak}}(\btheta) &\leq T_{\max},\\
0 \leq \mathrm{SOC}(t) &\leq 0.80 .
\end{align}
Protocols that violate hard operational constraints receive penalty objective
values and are not allowed to dominate feasible protocols.

% ============================================================================
% Section III: Electrochemical-Thermal-Aging Model
% ============================================================================
\section{Electrochemical-Thermal-Aging Model}
\label{sec:model}

% 3.1 Electrochemical Model
\subsection{Electrochemical Model}

The simulator uses the Single Particle Model with electrolyte (SPMe), a
reduced-order electrochemical model that preserves electrolyte concentration
dynamics while approximating each electrode by a representative spherical
particle~\cite{ref6}. SPMe offers a useful trade-off between physical fidelity
and computational cost for optimization. The implementation is based on
PyBaMM~\cite{ref7}. The Chen2020 parameter set is used for the LG
INR21700-M50 cell, and the Ecker2015 parameter set is used as an additional
cross-parameter validation case~\cite{ref8,ref9}.

% 3.2 Thermal Model
\subsection{Thermal Model}

A lumped thermal model is coupled to the electrochemical model:
\begin{equation}
\rho C_p \frac{\partial T}{\partial t}
=q_{\mathrm{gen}}-h_c(T-T_{\mathrm{air}}),
\label{eq:thermal}
\end{equation}
where $q_{\mathrm{gen}}$ is heat generation from electrochemical reactions and
ohmic loss, $\rho$ is density, $C_p$ is heat capacity, and $h_c$ is the
convective coefficient. The peak temperature rise objective is extracted as
$\Delta T_p=\max_t T(t)-T_{\mathrm{init}}$.

% 3.3 Aging Model
\subsection{Aging Model}

The aging objective is represented as a capacity-fade percentage:
\begin{equation}
Q_s=\frac{Q_{\mathrm{eff}}}
{\mathrm{Cap}(\bar{s},\bar{T},\bar{I})}\times 100\%,
\label{eq:aging}
\end{equation}
where $\bar{s}$, $\bar{T}$, and $\bar{I}$ are the mean SOC, temperature, and
current during charging. The effective capacity function follows an empirical
temperature- and current-dependent form used for fast-charging optimization.
This objective is not intended to replace long-horizon cycling experiments; it
provides a consistent simulation-side degradation signal for comparing
protocols under the same model assumptions.

% 3.4 Simulation Implementation
\subsection{Simulation Implementation}

For a protocol $\btheta$, the simulator runs three sequential stages. First, it
computes the SOC breakpoints from $\Delta s_1$, $\Delta s_2$, and the remaining
third-stage width. Second, it solves the SPMe-thermal model stage by stage with
constant currents $I_1$, $I_2$, and $I_3$. Third, it extracts $t_c$,
$\Delta T_p$, and $Q_s$ from the concatenated trajectory. The optimizer sees
the simulator as a black-box vector function, while feasibility flags and
penalty values prevent failed simulations from breaking the BO loop.

% ============================================================================
% Section IV: Proposed LLMBO-MO Framework
% ============================================================================
\section{Proposed LLMBO-MO Framework}
\label{sec:method}

% 4.1 Framework Overview
\subsection{Motivation and Framework Overview}

LLMBO-MO keeps the numerical BO loop as the decision authority and uses the LLM
only as a source of bounded domain knowledge. The closed loop consists of:
\begin{enumerate}
\item LLM warmstart generation, validation, repair, and portfolio selection.
\item Initial simulator evaluation and insertion into the observation database
$\bD$.
\item Riesz-weight sampling and augmented Tchebycheff scalarization.
\item GP fitting on scalarized observations.
\item LLM point or region guidance based on the current BO state.
\item Guarded region-lifted GP coupling or fallback to plain EI.
\item Simulator evaluation of the selected protocol and Pareto/HV update.
\end{enumerate}
Thus, the LLM influences initialization and local search preference, but the
final sample is still selected by a GP/EI acquisition calculation and verified
by the simulator.

% 4.2 Decomposition Strategy
\subsection{Riesz s-Energy Weight Generation and Tchebycheff Scalarization}

The objectives are first transformed to improve numerical conditioning:
\begin{equation}
\tilde f_1=\log_{10}(t_c),\quad
\tilde f_2=\Delta T_p,\quad
\tilde f_3=\log_{10}(Q_s).
\label{eq:log}
\end{equation}
The transformed objectives are normalized by the gap to an ideal point:
\begin{equation}
\bar f_i(\btheta)=
\frac{|\tilde f_i(\btheta)-z_i^\ast|}
{y_{\max,i}-y_{\min,i}},\quad i=1,2,3.
\label{eq:norm}
\end{equation}
Given a weight vector $\bw$ on the simplex, the scalarized target is
\begin{equation}
g^{\mathrm{tch}}(\btheta|\bw)
=\max_i\{w_i\bar f_i(\btheta)\}
\eta\sum_{i=1}^{3}w_i\bar f_i(\btheta),
\label{eq:tch}
\end{equation}
where $\eta=0.05$. The weight set $\bW$ is generated by minimizing the Riesz
$s$-energy on the simplex,
\begin{equation}
\min_{\bw^{(1)},\ldots,\bw^{(N)}}\sum_{i<j}
\|\bw^{(i)}-\bw^{(j)}\|^{-s},
\label{eq:riesz}
\end{equation}
which encourages even coverage of different Pareto trade-off directions.

% 4.3 GP Surrogate Model
\subsection{Gaussian Process Surrogate Model}

For each scalarized BO subproblem, LLMBO-MO fits a single-output GP on
$\{(\btheta_j,g_j^{\mathrm{tch}})\}_{j=1}^{|\bD|}$. The kernel is a Matern
$5/2$ kernel with automatic relevance determination:
\begin{equation}
k(\btheta,\btheta')=\sigma_f^2
\left(1+\sqrt{5}r+\frac{5r^2}{3}\right)\exp(-\sqrt{5}r),
\label{eq:matern}
\end{equation}
where $r$ is the normalized ARD distance. Hyperparameters are estimated by
maximizing the marginal likelihood with multiple restarts. The posterior mean
and variance are then used by EI and by the LLM-guidance coupling.

% 4.4 LLM Integration (Two Touchpoints)
\subsection{LLM Integration: Two-Touchpoint Architecture}

\subsubsection{Touchpoint 1: Warmstart Candidate Generation}
The warmstart prompt provides cell information, variable bounds, the dSOC
constraint, and electrochemical design principles such as decreasing current
profiles and safety margins. The LLM returns a pool of candidate protocols.
The system parses, clips, repairs, filters, and deduplicates these candidates.
A greedy portfolio selector then chooses $n_{\mathrm{ws}}$ candidates that are
both feasible and diverse:
\begin{equation}
J(\btheta)=
-\omega_d d_{\min}(\btheta,\mathcal{C}_{\mathrm{sel}})
\omega_p p_{\mathrm{soft}}(\btheta)
\omega_m p_{\mathrm{mono}}(\btheta).
\label{eq:portfolio}
\end{equation}
Here $d_{\min}$ rewards distance from already selected points,
$p_{\mathrm{soft}}$ penalizes soft dSOC violations, and $p_{\mathrm{mono}}$
penalizes non-monotone currents. If the LLM returns too few valid candidates,
hand-crafted physics heuristics and random/LHS points fill the initial design.

\subsubsection{Touchpoint 2: Iteration-Level Guidance}
At each BO iteration, the LLM receives the current weight vector, the best and
recent observations, Pareto-front context, and uncertainty hotspots. It returns
either a point preference $\btheta_{\mathrm{pref}}$ or a bounded region
$(\mathbf{lb},\mathbf{ub})$, along with a confidence score. Guidance must pass
parsing, confidence, boundary, width, and feasibility checks before it can
affect the acquisition function. Otherwise, the optimizer uses plain EI.

% 4.5 GP-LLM Coupling
\subsection{GP-LLM Coupling via Acquisition-Time Mean Shift}

Valid LLM guidance is lifted into the GP acquisition space through a local grid
$\bG$ around the preferred point or inside the preferred region. The coupling
strength is
\begin{equation}
\lambda=\mathrm{clip}\left(
\frac{c\kappa_c}{\sqrt{\mathbf{v}^{\top}\Sigma_{\bG\bG}\mathbf{v}}}
\rho^t,\lambda_{\min},\lambda_{\max}\right),
\label{eq:lambda}
\end{equation}
where $c$ is LLM confidence, $\mathbf{v}$ contains grid weights, and $\rho^t$
anneals LLM influence over time. For a candidate $\btheta$, the acquisition-time
mean is shifted as
\begin{equation}
\mu_c(\btheta)=\mu(\btheta)
-\lambda g_{\mathrm{gate}}m_{\mathrm{local}}(\btheta)q(\btheta)\hat\sigma_y .
\label{eq:shift}
\end{equation}
Only the mean used by the acquisition is modified; the GP posterior variance is
left unchanged. This design lets the LLM bias the candidate ranking without
pretending to have observed new data.

% 4.6 Acquisition Function
\subsection{Acquisition Function with Stagnation-Aware Exploration}

Expected improvement is computed from the acquisition-time mean:
\begin{equation}
\alpha_{\mathrm{EI}}(\btheta)
=(f_{\min}-\mu_c(\btheta))\Phi(z)+\sigma(\btheta)\phi(z),
\label{eq:ei}
\end{equation}
where $z=(f_{\min}-\mu_c(\btheta))/\sigma(\btheta)$. The final score combines
normalized EI with a guidance bonus and a risk penalty:
\begin{equation}
\alpha(\btheta)=\widehat{\alpha}_{\mathrm{EI}}(\btheta)
\beta_{\mathrm{guid}}(\btheta)-\beta_{\mathrm{risk}}(\btheta).
\label{eq:acq}
\end{equation}
Risk terms discourage candidates close to hard dSOC limits or with physically
implausible increasing current profiles. If hypervolume or scalarized
improvement stagnates, the local sampling width is expanded to encourage
exploration.

% 4.7 Overall Algorithm
\subsection{Overall Algorithm}

\begin{algorithm}[t]
\caption{LLMBO-MO}
\label{alg:llmbo}
\begin{algorithmic}[1]
\Require Simulator $\bS$, LLM $\bL$, budget $T$, warmstart count
$n_{\mathrm{ws}}$, weight set $\bW$
\Ensure Pareto front $\bP$ and hypervolume trace
\State Generate or load Riesz weight set $\bW$
\State Query LLM for warmstart pool and select $n_{\mathrm{ws}}$ candidates
\For{each selected warmstart protocol}
    \State Evaluate simulator and add observation to $\bD$
\EndFor
\For{$t=1$ to $T$}
    \State Select weight $\bw^{(t)}$ and compute scalarized targets
    \State Fit GP on the scalarized database
    \State Query LLM for point or region guidance
    \If{guidance passes all guards}
        \State Build region-lifted mean shift and acquisition prior
    \Else
        \State Use plain EI fallback
    \EndIf
    \State Select next protocol by maximizing acquisition score
    \State Evaluate simulator, update $\bD$, $\bP$, and HV
\EndFor
\State \Return $\bP$
\end{algorithmic}
\end{algorithm}

% ============================================================================
% Section V: Experiments
% ============================================================================
\section{Experiments}
\label{sec:experiments}

% 5.1 Experimental Setup
\subsection{Experimental Setup}

Experiments are conducted on SPMe-based simulators using the Chen2020 and
Ecker2015 parameter sets. Each BO run uses 56 simulator evaluations unless
otherwise stated, corresponding to six initialization evaluations and 50 BO
iterations. The primary metric is canonical hypervolume (HV), computed from the
nondominated set in the transformed objective space. Higher HV indicates a
larger dominated objective volume relative to the reference point.

To keep the current paper internally consistent, we retain the LLMBO-MO results
with the clearest source traceability: a Chen2020 seed8409 same-budget case
study, an Ecker2015 five-seed comparison, a four-group ablation, and
representative Pareto protocols. Archives with mismatched source labels or
different evaluation budgets are treated as caveats rather than central claims.

% 5.2 Baseline Methods
\subsection{Baseline Methods}

The principal BO baseline is ParEGO, which uses scalarization and acquisition
optimization without LLM knowledge. NSGA-II is included as an evolutionary
reference. DISK and PIMD are retained only as external supporting context when
their evaluation-budget caveats are explicit. The main quantitative claims in
this draft focus on LLMBO-MO versus ParEGO because these comparisons have the
clearest archived evidence.

% 5.3 HV Convergence Comparison
\subsection{HV Convergence Comparison}

Table~\ref{tab:chen_case} reports the retained Chen2020 advantage case. Under
the same seed8409 and 56-evaluation budget, the tuned GPT-4.1-mini LLMBO-MO run
achieves a higher canonical HV than the ParEGO reference run.

\begin{table}[htbp]
\caption{Chen2020 Seed8409 Case Study, 56 Evaluations}
\label{tab:chen_case}
\centering
\begin{tabular}{lcc}
\toprule
Method & Canonical HV & Improvement \\
\midrule
ParEGO reference & 0.3523111 & -- \\
LLMBO-MO & \textbf{0.3848256} & \textbf{+0.0325145} \\
\bottomrule
\end{tabular}
\end{table}

This result is used as a representative Chen2020 case study. It is not
presented as a claim that every Chen2020 five-seed aggregate favors LLMBO-MO,
because the current archive contains multiple result sets with different LLM
backends and source labels.

Table~\ref{tab:ecker} reports the strongest retained multi-seed result. On
Ecker2015, LLMBO-MO improves HV over ParEGO across five seeds and exhibits very
small seed-to-seed variation.

\begin{table}[htbp]
\caption{Ecker2015 Final Canonical HV, 5 Seeds, 56 Evaluations}
\label{tab:ecker}
\centering
\begin{tabular}{lccc}
\toprule
Method & Mean HV & Std HV & Pareto Size \\
\midrule
ParEGO & 1.5866 & 0.0116 & 34.6 \\
LLMBO-MO & \textbf{1.8684} & \textbf{0.0024} & 26.8 \\
\bottomrule
\end{tabular}
\end{table}

The Ecker2015 result gives a mean canonical-HV improvement of approximately
$+0.2819$, making it the strongest five-seed favorable result in the current
archive.

% 5.4 Pareto Front Quality
\subsection{Pareto Front Quality}

Table~\ref{tab:protocols} lists representative protocols from the retained
Chen2020 seed8409 case. The points are selected from fast, balanced, and
conservative regions of the Pareto front and are reported as
$(t_c,\Delta T_p,Q_s)$.

\begin{table*}[htbp]
\caption{Representative Chen2020 Pareto Protocols from the Seed8409 Case}
\label{tab:protocols}
\centering
\begin{tabular}{llccc}
\toprule
Method & Region & Charging time (s) & Peak temp. rise (K) & Capacity fade (\%) \\
\midrule
\multirow{3}{*}{LLMBO-MO}
& Fast & 2880 & 7.567 & 1.261 \\
& Balanced & 6112 & 2.857 & 0.571 \\
& Conservative & 7200 & 1.529 & 0.640 \\
\midrule
\multirow{3}{*}{ParEGO reference}
& Fast & 3304 & 6.315 & 1.024 \\
& Balanced & 5290 & 3.202 & 0.621 \\
& Conservative & 7109 & 1.537 & 0.651 \\
\bottomrule
\end{tabular}
\end{table*}

The balanced and conservative representatives show that LLMBO-MO can offer
competitive temperature-degradation trade-offs, not only a better scalar HV.
The fast representative is more aggressive and illustrates the expected
speed-aging trade-off at the fast-charging edge.

% 5.5 Ablation Study
\subsection{Ablation Study}

Table~\ref{tab:ablation} summarizes the four-group ablation. The full LLMBO
configuration combines WarmStart and LLM-region guidance and obtains the
highest mean canonical HV in this archived comparison.

\begin{table}[htbp]
\caption{Chen2020 Four-Group Ablation, 5 Seeds, 50 Iterations}
\label{tab:ablation}
\centering
\begin{tabular}{lccc}
\toprule
Variant & Mean HV & $\Delta$HV & Wins \\
\midrule
Baseline & 0.383635 & -- & -- \\
WarmStart & 0.390242 & +0.006607 & 4/5 \\
LLM\_Region & 0.386211 & +0.002576 & 4/5 \\
Full LLMBO & \textbf{0.393196} & \textbf{+0.009561} & 3/5 \\
\bottomrule
\end{tabular}
\end{table}

The ablation suggests that WarmStart is the most stable individual contributor,
while LLM-region guidance also provides positive mean improvement. Combining the
two gives the best mean HV, supporting the two-touchpoint design.

% 5.6 Computational Efficiency
\subsection{Computational Efficiency}

Runtime measurements on Chen2020 with a DeepSeek-V3 backend show that LLMBO-MO
requires additional wall-clock time because it includes LLM calls, guidance
parsing, region processing, and a richer acquisition pool. Table~\ref{tab:time}
reports these values as computational context only.

\begin{table}[htbp]
\caption{Runtime Context on Chen2020, 5 Seeds}
\label{tab:time}
\centering
\begin{tabular}{lcc}
\toprule
Method & Runtime mean (s) & Runtime std (s) \\
\midrule
NSGA-II & 194.3 & 11.8 \\
ParEGO & 252.4 & 17.7 \\
LLMBO-MO & 440.3 & 31.3 \\
\bottomrule
\end{tabular}
\end{table}

This overhead is a practical limitation of LLM-guided BO. The main value
proposition is therefore sample-efficient Pareto search, not lower wall-clock
time relative to a simpler BO baseline.

% ============================================================================
% Section VI: Conclusion
% ============================================================================
\section{Conclusion}
\label{sec:conclusion}

This paper presented LLMBO-MO, a large language model-augmented
multi-objective Bayesian optimization framework for battery fast-charging
protocol design. The framework uses the LLM at two guarded touchpoints:
warmstart generation and iteration-level point/region guidance. The optimizer
keeps the final decision inside a conventional BO loop through Riesz/Tchebycheff
scalarization, GP modeling, EI acquisition, risk penalties, and simulator
feedback. Retained experiments show favorable LLMBO-MO behavior in a Chen2020
seed8409 case study, strong five-seed improvement on Ecker2015, and an ablation
where the full two-touchpoint configuration achieves the highest mean HV.

Future work will add stronger statistical validation across more cell models,
study locally hosted LLM backends to reduce runtime overhead, and extend the
framework to variable-stage protocols with mixed discrete-continuous decision
spaces.

% ============================================================================
% References
% ============================================================================
\begin{thebibliography}{14}

\bibitem{ref1}
N.~Nitta \textit{et al.}, ``Li-ion battery materials: present and future,''
\textit{Materials Today}, vol.~18, no.~5, pp.~252--264, 2015.

\bibitem{ref2}
A.~Tomaszewska \textit{et al.}, ``Lithium-ion battery fast charging: A review,''
\textit{eTransportation}, vol.~1, p.~100011, 2019.

\bibitem{ref3}
B.~Shahriari \textit{et al.}, ``Taking the human out of the loop: A review of
Bayesian optimization,'' \textit{Proc. IEEE}, vol.~104, no.~1, pp.~148--175,
2016.

\bibitem{ref4}
J.~Knowles, ``ParEGO: A hybrid algorithm with on-line landscape approximation
for expensive multiobjective optimization problems,'' \textit{IEEE Trans. Evol.
Comput.}, vol.~10, no.~1, pp.~50--66, 2006.

\bibitem{ref5}
\usepackage[T1]{fontenc}
K.~Deb \textit{et al.}, ``A fast and elitist multiobjective genetic algorithm:
NSGA-II,'' \textit{IEEE Trans. Evol. Comput.}, vol.~6, no.~2, pp.~182--197,
2002.

\bibitem{ref6}
S.~G.~Marquis \textit{et al.}, ``An asymptotic derivation of a single particle
model with electrolyte,'' \textit{J. Electrochem. Soc.}, vol.~166, no.~15,
pp.~A3693--A3706, 2019.

\bibitem{ref7}
V.~Sulzer \textit{et al.}, ``Python Battery Mathematical Modelling (PyBaMM),''
\textit{J. Open Res. Softw.}, vol.~9, no.~1, p.~14, 2021.

\bibitem{ref8}
C.~H.~Chen \textit{et al.}, ``Development of experimental techniques for
parameterization of multi-scale lithium-ion battery models,''
\textit{J. Electrochem. Soc.}, vol.~167, no.~8, p.~080534, 2020.

\bibitem{ref9}
M.~Ecker \textit{et al.}, ``Parameterization of a physico-chemical model of a
lithium-ion battery: I. Determination of parameters from electrochemical
impedance spectroscopy,'' \textit{J. Power Sources}, vol.~295, pp.~145--154,
2015.

\bibitem{ref10}
T.~Brown \textit{et al.}, ``Language models are few-shot learners,'' in
\textit{Advances in Neural Information Processing Systems}, vol.~33, 2020.

\bibitem{ref11}
X.~Li \textit{et al.}, ``Data-driven multi-objective optimization for battery
charging,'' \textit{IEEE Trans. Intell. Veh.}, vol.~8, no.~3,
pp.~2456--2467, 2023.

\bibitem{ref12}
Y.~Zhang \textit{et al.}, ``Physics-informed model-driven optimization for fast
charging protocol design,'' \textit{Appl. Energy}, vol.~335, p.~120728, 2023.

\bibitem{ref13}
J.~Jiang \textit{et al.}, ``Bayesian optimization for charging protocol
design,'' \textit{IEEE Trans. Ind. Electron.}, vol.~71, no.~4,
pp.~4120--4131, 2024.

\bibitem{ref14}
S.~Krishnamoorthy \textit{et al.}, ``Large language models for automated design
optimization,'' \textit{arXiv preprint arXiv:2402.06196}, 2024.

\end{thebibliography}

\end{document}
